Surrogate-assisted damage-threshold estimation is investigated for ultrashort laser interactions with metallic thin films. The physical formulation is based on the parabolic two-temperature model (pTTM) for axisymmetric ultrashort-pulse laser heating, reduced to one- and two-dimensional computational problems. It describes nonequilibrium electron--lattice energy exchange after femtosecond laser deposition and combines temperature-dependent optical and thermophysical response with ballistic and finite-thickness source corrections.

Stable solvers are developed for linear and nonlinear pTTM threshold calculations. The linear solver is a two-field monolithic implicit Crank--Nicolson discretisation with Picard iteration and Anderson acceleration. The nonlinear solver is a two-field sequential IMEX Heun--Crank--Nicolson scheme with a strong Gauss--Seidel predictor, treating diffusion and electron--phonon coupling implicitly while retaining the bounded source contribution explicitly. Stiffness analysis motivates this split. Refinement studies show observed second-order behaviour in time and space. For the tested nonlinear configuration, the selected split is about \(1.15\times\) as fast as the fully implicit Crank--Nicolson alternative at comparable accuracy.

Coupled with exponential bracketing and bisection, the solvers provide melting-onset threshold estimates for representative metals. The surrogate models learn the maximum lattice temperature \(T_{l,\max}\) over synthetic two-dimensional pTTM datasets, and threshold fluences are recovered by inversion. The linear study uses LightGBM and a residual multilayer perceptron. A nonlinear extension is also included at the \(\num{20000}\) accepted-sample scale, using temperature-dependent \(C_e(T_e)\) and \(G(T_e)\), Lorentz--Drude optical descriptors, source-term-informed feature engineering, LightGBM, and a standard multilayer perceptron. On independent test splits, the linear residual MLP and LightGBM obtain \(R^2=0.997\) and \(0.990\), while the nonlinear MLP and LightGBM obtain \(R^2=0.982\) and \(0.966\), respectively.

The linear surrogate is accurate within the retained synthetic target policy but exhibits high-temperature target-tail censoring and material-support limitations for several transition-metal validation curves. The nonlinear analysis improves the interpretation of Ni and Ti relative to the linear failure case and reports internally generated threshold estimates for representative nonlinear material cases.
